% 正赛文稿：正方。环节标题按 recruit-3v3.md「各持方文稿标题」原样使用。
\documentclass[script,pro]{debate}
\motion{“磕CP”热潮折射出人们对爱情的更深信仰 / 更深怀疑}
\stance{正方}{更深信仰}
\meta{赛制}{招新赛 3v3}
\meta{口径来源}{备赛文档 第十节 正方口径表}
\begin{document}
\debatetitle
\begin{thesisbox}[全场一句话立论]人们嗑的是最严格的爱情标准，疑的只是遇不遇得到；疑在上面，信在底下。\end{thesisbox}

\begin{speech}{1. 正一开篇立论稿}{正文 823 字 / 180 秒}
谢谢主席，各位好。

这道题问的是，嗑CP这股热潮底下，藏的是对爱情更深的信，还是更深的疑。先把几个词说清楚。嗑CP，是以旁观者的身份，给虚构角色或者真实人物配对，想象并期待他们之间的浪漫关系，并从中得到情感满足。对爱情的信仰，是相信真正的爱情存在，值得期待和投入。怀疑正好反过来：不相信它存在，或者不相信它靠得住、值得投入。

我方承认，热潮里两种心态都有。所以“更深”，比的是哪一种在底下。判断标准是：哪一种心态更能解释热潮的核心特征，也就是人们嗑的是什么样的爱，为什么用旁观的方式去爱，并且它是另一种心态的来源。在这个标准下，我方要证明，嗑CP的动力来自相信真爱存在，犹豫疑的是遇见，不是爱情本身；对方要证明，嗑CP的动力来自不信真爱靠得住。

我方第一点，嗑的是理想。

嗑谁，嗑什么，是每个人自己挑的。挑什么，就暴露了心里的标准。你想想，大家嗑得最上头的是哪一种？是认定了就不换人的，是兜兜转转还能走到一起的，是两个人都在往对方那边走的。心理学家 Sprecher 和 Metts 在一九八九年编过一份浪漫信念量表，四条：唯一，爱能找到出路，把对方理想化，一见钟情。一部让人上头的CP剧，就是把这张量表演了一遍。人民大学的许冠文在二〇二三年深访了五十三位女性CP粉，他说，CP更像是一个把她们的期待具象化的载体。

“双向奔赴”，是《咬文嚼字》评出的二〇二三年十大流行语，释义是“相互爱慕、相互亲近的美好愿望”。标准没有降，是\stress{升了}。

我方第二点，疑的是条件。

热潮里的犹豫是真的，我方不否认。但要问一句：疑的是什么？复旦发展研究院二〇二四年的调查，八千多名单身青年里，将近八成说想谈恋爱。问为什么没脱单，排第一的是社交圈子单一，六成多的人选了；后面是很难遇到三观契合的对象，不太擅长和异性交流。没有一条是爱情是假的。

中国青年报采访过一位二十八岁的姑娘，化名张钰。负面的婚恋新闻常上热搜，让她对感情提不起信心。可她最后说的是：想遇到对的人一起走，两人彼此信任。说白了，卡住她的是门口，不是门里没有人。疑在上面，信在底下。

我方全场都请对方回答一个问题：如果人们不信爱情，为什么偏偏嗑“唯一”和“双向奔赴”，而不嗑一对互相算计的人？

这一代人可以晚一点恋爱，可以谨慎一点。可他们还在为别人的好结局高兴。\stress{这份高兴，就是信}。谢谢。
\end{speech}
\begin{contingency}
\ifthen{反方把“信仰”定义成“必须自己去谈”}{信仰不以亲身得到为前提；按这个标准，所有单身的人都不信爱情，这是拿婚恋率换掉了爱情观。请回到中立判准：看热潮的特征。}
\ifthen{反方用婚恋率、结婚登记说话}{结婚登记跟着条件走，二〇二五年全国通办以后一年多了六十五万多对；爱情观不会一年变一次。}
\ifthen{反方拿出我方没见过的数据}{先问机构、年份、问的是不是嗑CP的人；再问这个数测的是想不想，还是信不信。}
\ifthen{论点一被打成“想象不用兑现”}{收缩到论点二：疑的对象是遇见，用判准说明底下是信。}
\end{contingency}

\begin{stage}{2. 正一被质询防守预案}{反二质询，90 秒双边计时}
\textbf{原则}：先答是或否，再用一句话守住前提；不拖时间，被打断就停。
\begin{qa}
\qarow{嗑CP的人，自己进入那段关系吗？}{不进入。看的人不上场，不等于不信；我方看的是他们选了看什么。}
\qarow{那您方承认嗑CP是旁观？}{承认，定义里就写了旁观者。旁观是形式，我方问的是形式底下的心态。}
\qarow{将近八成想谈恋爱，想谈就是信吗？}{想谈加上不信它是假的，就是信；那份调查里没脱单的理由，没有一条是爱情是假的。}
\qarow{量表测的是人对自己的信念，不是对CP的？}{对，量表说明的是浪漫信念有哪几条；CP剧演的正是这几条，这是对照，不是测量。}
\qarow{“双向奔赴”后来也用在别的地方，跟爱情有关吗？}{释义第一句就是多用于人与人之间，表达相互爱慕的愿望。用得越广，越说明这个理想被接受。}
\qarow{张钰对感情提不起信心，不就是怀疑？}{她疑的是热搜里那些婚姻，也就是她能不能遇到对的人；她最后说的是想遇到对的人。}
\end{qa}
\end{stage}

\begin{stage}{3. 正二质询反一问题链}{90 秒，双边计时}
\textbf{总目标}：拿到两个承认：大家嗑的是专一、双向的爱；疑的对象是遇见，不是爱本身。问题短，对方不答时先复述一次再礼貌换问。
\begin{chain}{链一：内容链}{让对方承认嗑的是理想}
\q{大家嗑得最多的，是专一的一对，还是互相算计的一对？}{答专一：下一问；答都有：那热搜上最火的是哪种？}
\q{一个不信爱情的人，会为专一的一对上头吗？}{答会：那他上头的是专一本身，您方承认他在乎专一；答不会：谢谢，收束}
\q{在乎专一，算不算相信专一的爱值得？}{答算：记下承认；答不算：请评委记住对方把在乎和相信分开了}
\close{感谢，对方承认热潮里被嗑的是专一的爱。}
\end{chain}
\begin{chain}{链二：对象链}{让对方承认疑的是遇见}
\q{没脱单的原因里，排第一的是社交圈子单一，对吗？}{答对：下一问；答不知道：复旦发展研究院二〇二四年，六成多}
\q{圈子窄，是怀疑爱情，还是怀疑遇不遇得到？}{答遇不遇得到：收束；答怀疑爱情：圈子变宽，爱情就变真了吗？}
\q{看球的人不上场，能说他不信足球吗？}{答不能：旁观不等于怀疑；答能：那全场观众都不信足球}
\cut{好，我换个问法。}
\close{感谢，对方承认疑的是遇见，旁观不等于怀疑。}
\end{chain}
\end{stage}

\begin{stage}{4. 正三对辩要点}{双方各 90 秒}
\textbf{任务}：开全场第一个战场：\textbf{内容与形式，谁更说明心态}。每次发言先回应一句，再抛一个问题。
\begin{points}{进攻点（4–5 个，每个 15–20 秒）}
\item 大家嗑的是专一、双向、好结局。不信爱情的人，为什么偏偏嗑这些？
\item 旁观是形式。看球的人不上场，不说明他不信足球；请对方说出旁观和怀疑之间缺的那一步。
\item 没脱单的理由排前三的是圈子、遇见、交流，没有一条是爱情是假的。对方的怀疑，证据在哪？
\item 对方说不敢，不敢说明当真。一个不当真的东西，谁会怕它？
\item 更深比的是谁是谁的来源：是因为想要才嗑，还是因为不信才嗑？
\end{points}
\begin{points}{备用点（6–8 个）}
\item 对方说画好了终点：画好终点，是相信爱能走到终点。
\item 对方说代餐：人只给想要的东西找代餐。
\item 对方说婚恋率降了：结婚不等于爱情，二〇二五年全国通办后一年多了六十五万多对，数字跟着条件走。
\item 对方说“我又相信爱情了”的“又”：信能被唤回来，说明它一直在底下。
\item 对方说两百多份问卷里一半只嗑不谈：那是便利样本，也不是全体。
\item 对方说叶公好龙：叶公怕真龙；年轻人不怕爱，怕将就。
\item 对方说一成多不谈恋爱：那近九成呢？
\end{points}
\begin{points}{对方最可能问的三个问题 → 回应}
\item 你信爱情，为什么不去谈？ → 圈子窄、时间少，卡在遇见；嗑和谈不冲突。
\item 想象里标准高，不用兑现吧？ → 标准被带回现实：宁缺毋滥，拒绝将就就是在用这把尺子。
\item 许冠文说现实的亲密关系每一步都充满风险，不是疑吗？ → 同一句话的前半是“我们对这种良好的亲密关系赋予了很高的价值”；风险大，说明把爱当真。
\end{points}
\end{stage}

\begin{kit}{5. 正二申论}{套件 · 组装后 389–415 字 / 90 秒}
\assembly{开头 → 回应反二申论（任选 1） → 拆反方立论（任选 1） → 推进论点一 → 收尾}
\begin{fixed}{开头}
谢谢主席，各位。我这一篇做两件事：先看反方的立论撑不撑得住，再把我方第一点往前推一步。
\end{fixed}
\begin{pick}{回应反二申论}{1}
\begin{module}{如果反二申论主打“想象里的标准不用兑现”}
反二刚才说，想象里的标准不用兑现，所以不算信。我想问一句，信仰什么时候需要先兑现？没有见过，却还愿意相信，这恰恰就是信仰和经验的区别。先兑现再相信的，那叫验收。
\end{module}
\begin{module}{如果反二申论主打“想谈不等于相信”}
反二刚才说，想谈不等于相信。这一点我方同意，所以我方从来不只看想不想，我方看的是理由。没脱单的理由里，排前三的是圈子、遇见和交流，没有一条是爱情是假的。
\end{module}
\begin{fallback}{反二申论主打不在以上两条时}
反二这一篇，我方先记下来，放到判准底下看：它说明的是人们怎么嗑，还是人们为什么嗑？判准两样都问，而我方今天最关心的，是后一个问题。
\end{fallback}
\end{pick}
\begin{pick}{拆反方立论}{1}
\begin{module}{如果反方立论的主干是“旁观就是怀疑”}
反方的立论说，嗑CP是旁观，所以是怀疑。这里缺了一步。看球的人不上场，没有人说他不信足球。旁观只告诉我们人站在哪里，不告诉我们人信什么。要说旁观就是怀疑，反方得证明：这些人站在场外，是因为觉得场上的东西是假的。而不是因为不敢、嫌麻烦、圈子窄，这些都是门口的事。
\end{module}
\begin{module}{如果反方立论的主干是“怕风险就是怀疑”}
反方的立论说，人们怕风险，所以不信爱情靠得住。这一步也跳过去了。怕手术的人，不是不信医学。恰恰相反，觉得爱情风险大，是因为把它当真，当成一件会让人受伤、也会让人变好的事。一个人真觉得爱情是假的，他会说无所谓，不会说风险太大。说白了，怕是信的影子。
\end{module}
\begin{fallback}{反方立论的主干不在以上两条时}
反方的立论，请评委放回判准里看。判准问的是两件事：人们嗑的是什么样的爱，为什么用旁观的方式去爱。反方如果只讲了嗑的方式，没讲嗑的内容，那就只回答了一半。而另一半，大家嗑的是专一，是双向，是好结局，这一半在我方这边，请反方来碰一碰。
\end{fallback}
\end{pick}
\begin{fixed}{推进论点一}
我方第一点，嗑的是理想，今天再往前走一步：这套标准没有留在屏幕里，它被带回了现实。中国青年报社会调查中心二〇二四年的调查，问年轻人对一段恋爱有什么期待，将近六成说希望互相欣赏，差不多同样多的人说期待彼此包容。你想想，这不就是双向奔赴吗？中国青年报采访过的一位单身青年说，她对感情的态度是“不急不躁，宁缺毋滥”。宁缺毋滥，就是拿着尺子，拒绝将就。
\end{fixed}
\begin{fixed}{收尾}
所以请反方回答：如果不信爱情，为什么还要守着这把尺子？说到底，一个真不信的人，早就把尺子扔了。
\end{fixed}
\end{kit}

\begin{stage}{6. 正二被质询防守预案}{反三质询，90 秒双边计时}
\textbf{原则}：申论里说过的数字这一分钟就会被问，先答出处，再守前提。
\begin{qa}
\qarow{中青报那项调查，样本多少？}{文章没公开样本量。我方只用它说年轻人期待的方向，不用它说具体多大比例。}
\qarow{宁缺毋滥，不就是一直缺？}{缺是暂时的，标准是自己的。不将就，说明相信有对的人值得等。}
\qarow{期待变得更自信、更有安全感，这是爱情还是自我？}{两件事可以同时成立，爱一个人，同时变成更好的自己，本来就是双向。}
\qarow{您方承认有人因为嗑CP不谈恋爱吗？}{承认有，有人就有。题目比的是多数人底下的心态，不是有没有。}
\qarow{画好终点，是不是不敢面对不确定？}{画好终点，是相信爱能走到终点。不信的人，不会费心去画。}
\end{qa}
\end{stage}

\begin{stage}{7. 正三质询反二问题链}{90 秒，双边计时}
\textbf{总目标}：把“怕”变成“当真”，把人数拉回判准。
\begin{chain}{链一：当真链}{让对方承认怕风险说明把爱当真}
\q{您方说怀疑包括不信它靠得住，对吗？}{答对：下一问}
\q{怕受伤的人，是把爱当真，还是当假？}{答当真：下一问；答当假：那假的东西怎么会伤人？}
\q{当真的东西，能说他不信它存在吗？}{答不能：收束；答能：请评委记住，当真也被算成不信}
\close{感谢，对方承认怕的人是把爱当真的。}
\end{chain}
\begin{chain}{链二：样本链}{把数据拉回题目主体}
\q{一成多不谈恋爱，剩下的多数没说不谈，对吗？}{答对：下一问；答不对：请对方报出另一个数}
\q{这份调查问的是嗑CP的人吗？}{答不是：收束；答是：请给出处}
\close{感谢，对方的人数既不是多数，也不是嗑CP的人。}
\end{chain}
\end{stage}

\begin{stage}{8. 正一对辩要点}{双方各 90 秒}
\textbf{任务}：收拢前两段：把全场带回“更深”，也就是谁是谁的来源。
\begin{points}{进攻点（4–5 个，每个 15–20 秒）}
\item 更深比的是来源。是因为想要才嗑，还是因为不信才嗑？
\item 对方全场说的是不敢、麻烦、圈子窄，这些都在门口。门里有没有人，对方证明了吗？
\item 嗑CP的人嗑的是好结局。相信好结局存在，算信还是疑？
\item 我方第一个问题还在：不信爱情，为什么偏偏嗑“唯一”和“双向奔赴”？
\end{points}
\begin{points}{备用点（6–8 个）}
\item 对方说屏幕那头：屏幕那头的，是我们希望这头也有的样子。
\item 对方说保险：看电影想要好结局，不说明不信生活。
\item 对方说一半只嗑不谈：两百多份学生问卷，不能代表热潮。
\item 对方说“又”相信：能被唤回的信，一直在底下。
\item 对方说现实太险：险，是因为当真。
\item 对方说信仰要行动：拒绝将就也是行动。
\end{points}
\begin{points}{对方最可能问的三个问题 → 回应}
\item 信在底下，证据在哪？ → 嗑的内容和不脱单的理由：前者是理想，后者是门口。
\item 为什么不去谈？ → 卡在遇见；嗑和谈不冲突。
\item 屏幕那头的信，关掉屏幕还剩多少？ → 剩下的是那把尺子：宁缺毋滥。
\end{points}
\end{stage}

\begin{kit}{9. 正三申论}{套件 · 组装后 387–396 字 / 90 秒}
\assembly{开头 → 处理反方最强点（任选 1） → 推进论点二 → 收尾}
\begin{fixed}{开头}
谢谢主席，各位。我来处理反方到现在最有力的一点，然后把我方第二点往前推。
\end{fixed}
\begin{pick}{处理反方最强点}{1}
\begin{module}{如果反方全场最有力的是“画好终点”那句访谈}
反方最有力的，是那句“已经给他们画好了终点”。我方承认，这句话说的是想要确定。但请看画的是什么终点：不会解散，不会跑偏，走到最后。一个人如果觉得爱情终究会散，他画的终点就不会是这样。画好终点的人，是在替爱情写一个它本来就配得上的结局。想要确定，是因为现实里暂时没见到；想要的偏偏是这个结局，是因为相信它值得。
\end{module}
\begin{module}{如果反方全场最有力的是“想谈却不去谈”}
反方最有力的，是一句追问：既然想谈，为什么不去谈？我方的回答很简单，因为卡在遇见。圈子窄，时间少，遇不到三观合的人，这是现实的门槛，不是对爱情的判决。而且嗑CP和谈恋爱不是二选一，一个人完全可以一边守着标准，一边等着那个人。反方要证明的，不是有人没去谈，而是没去谈的原因，是他觉得爱情是假的。这一步，请反方给出来。
\end{module}
\begin{fallback}{反方最有力的一点不在以上两条时}
嗑CP的样子，隔着屏幕，结局写好，随时能退，这些我方从一开始就不否认。但判准还有另一半：人们在那块屏幕上嗑的是什么？是专一，是双向，是好结局。形式说明人站在哪里，内容说明人信什么。一个不信爱情的人，不会在屏幕前守着这样的标准，更不会为它熬夜。嗑的内容，才是这场热潮的心。这一半，请反方正面回答。
\end{fallback}
\end{pick}
\begin{fixed}{推进论点二}
我方第二点，疑的是条件，今天再补一层。中国科学院心理研究所二〇二四年调查了五万五千多名大学生，发现不愿意脱单的人，比不愿意结婚、不愿意生育的人都少。也就是说，婚、育、恋三样里，被放下得最少的就是爱情。人们放下的，是制度和成本，留下的是爱。这也是为什么嗑CP没有让人远离生活。两位传播学者 Horton 和 Wohl 在一九五六年就写过：对大多数观众，这种隔着屏幕的关系是日常生活的补充，日常关系里的预设，会在里面被重新确认。
\end{fixed}
\begin{fixed}{收尾}
说到底，全场的分歧已经很清楚：疑在上面，信在底下。反方要赢，就得证明一件事：不是门口太挤，而是门里真的没有人。
\end{fixed}
\end{kit}

\begin{stage}{10. 正三被质询防守预案}{反一质询，90 秒双边计时}
\textbf{原则}：先答是或否；三辩的数字都在申论里，出处背熟。
\begin{qa}
\qarow{中科院那份调查，问的是信不信爱情吗？}{问的是意愿。我方用它说明人们先放下的是婚育，爱情被留得最多。}
\qarow{Horton 和 Wohl 也说变成替代就是病态，对吗？}{对。所以他们把多数人的情况叫补充；反方要证明多数是替代。}
\qarow{一九五六年的电视研究，能用在今天吗？}{机制没变：隔着屏幕的关系，对多数人是补充。反方的立论也用了同一篇。}
\qarow{边嗑边等，数据在哪？}{我方没有这个比例。我方的数据是将近八成单身青年想谈，而理由卡在遇见。}
\qarow{不愿脱单的少，不就是说还是有人不愿？}{有。题目比的是底下是什么，不是有没有人犹豫。}
\end{qa}
\end{stage}

\begin{stage}{11. 正一质询反三问题链}{90 秒，双边计时}
\textbf{总目标}：全场最后一次质询，为结辩取两样东西：对方承认热潮里也有信；对方说不清谁是谁的来源。
\begin{chain}{链一：结局链}{取“相信好结局”的承认}
\q{嗑CP的人最想看到的结局，是好还是坏？}{答好：下一问；答都有：那热度最高的是哪种？}
\q{想看好结局，是相信它存在，还是不信？}{答相信：收束；答不信：不信的东西为什么想看？}
\close{感谢，对方承认嗑的人相信好结局存在。}
\end{chain}
\begin{chain}{链二：来源链}{逼对方说谁在底下}
\q{您方承认热潮里也有对爱情的信吗？}{答承认：下一问；答不承认：那“我又相信爱情了”是谁说的？}
\q{是因为不信才嗑，还是因为想要才嗑？}{答因为不信：那为什么嗑的是专一？答想要：收束}
\cut{好，我换个问法：先有想要，还是先有不敢？}
\close{感谢，对方承认信在热潮里；谁是来源，请评委在结辩里听。}
\end{chain}
\end{stage}

\begin{stage}{12. 正二对辩要点}{双方各 90 秒}
\textbf{任务}：结辩前最后一次交锋，守住判准的两半：内容和来源。不开新战场。
\begin{points}{进攻点（4–5 个，每个 15–20 秒）}
\item 判准有两半：嗑什么、为什么嗑。对方只讲了怎么嗑。
\item 对方的每一个理由，不敢、麻烦、圈子窄，都在门口。
\item 大家嗑的是专一和双向，不信爱情的人不会守这样的标准。
\item 宁缺毋滥，是在用尺子，不是扔掉尺子。
\end{points}
\begin{points}{备用点（6–8 个）}
\item 对方说旁观：看球的人不上场，不说明不信足球。
\item 对方说保险：怕受伤，是把爱当真。
\item 对方说代餐：人只给想要的东西找代餐。
\item 对方说人数：一成多不谈，不是多数；那份调查也不是问嗑CP的人。
\item 对方说“又”：能被唤回，说明一直在。
\item 对方说屏幕那头：那头的，是这头想要的样子。
\end{points}
\begin{points}{对方最可能问的三个问题 → 回应}
\item 信而不做，算什么信？ → 拒绝将就就是在做；信仰不以先兑现为条件。
\item 你方的信证据在哪？ → 嗑的内容是理想，不脱单的理由是门口。
\item 嗑CP是不是在逃避？ → 对多数人是补充，不是替代；逃避需要反方证明多数。
\end{points}
\end{stage}

\begin{kit}{13. 正一结辩}{套件 · 组装后 387–406 字 / 90 秒}
\assembly{归纳分歧（任选 1） → 处理最强点（任选 1） → 回应反一结辩（任选 1） → 临场 → 论点收束 → 价值升华 → 结尾}
\begin{pick}{归纳分歧}{1}
\begin{module}{如果全场主战场落在“旁观算不算怀疑”}
谢谢主席。今天吵得最久的，是旁观算不算怀疑。旁观只说明人站在哪里。站在场外的原因，是门口的事，不是门里的事。
\end{module}
\begin{module}{如果全场主战场落在“想谈却不去谈”}
谢谢主席。今天吵得最久的，是想谈却不去谈。不去的原因是圈子、时间、遇见，都是门口的事，不是对爱情的判决。
\end{module}
\begin{fallback}{主战场不在以上两条时}
谢谢主席。今天的分歧其实就一句：热潮底下，是想要，还是不敢。我方说，先有想要，才会有不敢。想要，在底下。
\end{fallback}
\end{pick}
\begin{pick}{处理最强点}{1}
\begin{module}{如果反方全场最有力的是“画好终点”}
反方最有力的是画好终点。可画的是什么终点？不散，不跑偏，走到最后。一个觉得爱情终究会散的人，画不出这样的终点。想要确定，是因为暂时没见到；想要这个结局，是因为信它值得。
\end{module}
\begin{module}{如果反方全场最有力的是“怕风险就是不信靠得住”}
反方最有力的，是怕风险。我方承认，这一代人怕受伤。可怕受伤的，是把爱当真的人。一个不当真的东西，没有人会怕它伤人。怕，是信投下的影子，影子越长，说明那束光越亮。
\end{module}
\begin{fallback}{反方最有力的一点不在以上两条时}
反方最有力的那一点，请评委放回判准的另一半去量。隔着屏幕，结局写好，这些我方从不否认。可人们嗑的内容，是专一，是双向，是好结局。内容说明人信什么。
\end{fallback}
\end{pick}
\begin{pick}{回应反一结辩}{1}
\begin{module}{如果反一结辩收在“信在别人身上”}
反一最后说，信在别人身上。可一个人会为别人的好结局高兴，是因为他相信好结局存在，也相信它值得。
\end{module}
\begin{module}{如果反一结辩收在“保险式的爱”}
反一最后说，这是一份保险式的爱。可保险保的是结局，说明人们认定这个结局值得保，值得等。
\end{module}
\begin{fallback}{反一结辩的收束不在以上两条时}
反一的结辩，请评委放回判准里听：它说明的是人们怎么嗑，还是人们为什么嗑。判准要的是后一个。
\end{fallback}
\end{pick}
\live{40}
\begin{fixed}{论点收束}
请评委看，我方两点到最后都还站着。第一点，嗑的是理想，标准没有降，是升了。第二点，疑的是条件，门口有人在等，门里没有人说爱是假的。所以这股热潮底下，是信。
\end{fixed}
\begin{fixed}{价值升华}
你想想这样一个很小的画面。晚上十一点，宿舍熄了灯，一个女生缩在被子里，把一段剪辑看了第三遍。那两个人最后在雨里撑开了同一把伞。她不会去谈一段将就的恋爱。可她也没有关掉屏幕，说一句都是假的。她只是在等。我们今天争的，不是年轻人谈了多少恋爱。是这个等的人，心里还有没有一把尺子。
\end{fixed}
\begin{fixed}{结尾}
说到底，她还在等，就说明她还相信。谢谢。
\end{fixed}
\end{kit}
\begin{contingency}
\ifthen{时间不够}{先把回应组压成一句，临场位只说一句；价值升华与结尾不删。}
\end{contingency}
\end{document}
